%%%%%%%%%%%%%%%%%%%%%%%% referenc.tex %%%%%%%%%%%%%%%%%%%%%%%%%%%%%%
% sample references
% %
% Use this file as a template for your own input.
%
%%%%%%%%%%%%%%%%%%%%%%%% Springer-Verlag %%%%%%%%%%%%%%%%%%%%%%%%%%
%
% BibTeX users please use
% \bibliographystyle{}
% \bibliography{}
%
%\biblstarthook{References may be \textit{cited} in the text either by number (preferred) or by author/year.\footnote{Make sure that all references from the list are cited in the text. Those not cited should be moved to a separate \textit{Further Reading} section or chapter.} The reference list should ideally be \textit{sorted} in alphabetical order -- even if reference numbers are used for the their citation in the text. If there are several works by the same author, the following order should be used: 
%\begin{enumerate}
%\item all works by the author alone, ordered chronologically by year of publication
%\item all works by the author with a coauthor, ordered alphabetically by coauthor
%\item all works by the author with several coauthors, ordered chronologically by year of publication.
%\end{enumerate}
%The \textit{styling} of references\footnote{Always use the standard abbreviation of a journal's name according to the ISSN \textit{List of Title Word Abbreviations}, see \url{http://www.issn.org/en/node/344}} depends on the subject of your book:
%\begin{itemize}
%\item The \textit{two} recommended styles for references in books on \textit{mathematical, physical, statistical and computer sciences} are depicted in ~\cite{science-contrib, science-online, science-mono, science-journal, science-DOI} and ~\cite{phys-online, phys-mono, phys-journal, phys-DOI, phys-contrib}.
%\item Examples of the most commonly used reference style in books on \textit{Psychology, Social Sciences} are~\cite{psysoc-mono, psysoc-online,psysoc-journal, psysoc-contrib, psysoc-DOI}.
%\item Examples for references in books on \textit{Humanities, Linguistics, Philosophy} are~\cite{humlinphil-journal, humlinphil-contrib, humlinphil-mono, humlinphil-online, humlinphil-DOI}.
%\item Examples of the basic Springer style used in publications on a wide range of subjects such as \textit{Computer Science, Economics, Engineering, Geosciences, Life Sciences, Medicine, Biomedicine} are ~\cite{basic-contrib, basic-online, basic-journal, basic-DOI, basic-mono}. 
%\end{itemize}
%}

\begin{thebibliography}{99}

\bibitem{Baier} T.~Baier, C.~Florentino, J.~Mourao, J.~Nunes, {\em Toric K\"ahler metrics seen from infinity, quantization and compact tropical amoebas}. J. Differential Geom. 89 (2011), no. 3, 411--454.
 
 
\bibitem{BN} M.~Baker, S.~Norine, {\em Riemann-Roch and Abel-Jacobi theory on a finite graph},  Advances in Mathematics 215 (2007), 766--788.

\bibitem{BS} M.~Baker, F.~Shokrieh, {\em Chip-firing games, potential theory on graphs, and spanning trees.} 
J. Combin. Theory Ser. A 120 (2013), no. 1, 164--182. 

\bibitem{Besicovitch}
A.S. Besicovitch, {\em Almost Periodic Functions,} Cambridge University Press, 1932.

\bibitem{Shor} A.~Bjorner, L.~Lovasz,  P.~W.~Shor, {\em Chip-firing games on graphs}, European
J. Combin.,  12(4), (1991), 283--291.

\bibitem{bohr}
H. Bohr, {\em Almost Periodic Functions,} Chelsea Pub. Co., New York, 1947.

\bibitem{Bochner} S.~Bochner,  {\em Beitrage zur Theorie der fastperiodischen Funktionen}, Math. Annalen, 96: 119--147.

\bibitem{B2} E.~Bombieri, {\em The classical theory of Zeta and L-functions},
Milan J. Math. Vol. 78 (2010) 11-59.

\bibitem{B3} E.~Bombieri, J.~Lagarias {\em Complements to Li's criterion for the Riemann hypothesis}. J. Number Theory 77 (1999), no. 2, 274-287.


\bibitem{BJ} V.~Borchsenius, B.~Jessen, { \em Mean motion and values of the Riemann zeta function}. Acta Math. 80 (1948), 97--166.

\bibitem{BC} J.B.~Bost, A.~Connes, {\em Hecke algebras, Type III
factors and phase transitions with spontaneous symmetry breaking in
number theory}, Selecta Math. (New Series) Vol.1 (1995) N.3,
411--457.
%\bibitem{Cint} A.~Connes, {\em Sur la th\'eorie non commutative de l'int\'egration}
%Alg\`ebres d'op\'erateurs (S\' em., Les Plans-sur-Bex, 1978), p. 19-143, Lecture Notes in %Math., 725, Springer, Berlin, 1979.

\bibitem{Co-zeta} A.~Connes, {\em Trace formula in noncommutative
geometry and the zeros of the Riemann zeta function}.  Selecta Math.
(N.S.)  5  (1999),  no. 1, 29--106.

\bibitem{cmindex}  A. Connes, H. Moscovici, {\em
 The local index formula in noncommutative
geometry},  GAFA, Vol. 5 (1995), 174--243.


\bibitem{CoMo1} A.~Connes, H.~Moscovici,
{\em Modular Hecke Algebras and their Hopf Symmetry},  Mosc. Math.
J.  4  (2004),  no. 1, 67--109, 310.

\bibitem{CoMo2} A.~Connes, H.~Moscovici,
{\em Rankin-Cohen Brackets and the Hopf Algebra of Transverse
Geometry},  Mosc. Math. J.  4  (2004),  no. 1, 111--130, 311.

\bibitem{CoMo3} A.~Connes, H.~Moscovici,
{\em Transgressions of the Godbillon-Vey class and Rademacher
functions}, in "Noncommutative Geometry and Number Theory",
pp.79--107. Vieweg Verlag, 2006.

\bibitem{cmbook}  A.~Connes, M.~Marcolli, {\em Noncommutative Geometry, Quantum Fields and Motives} American Mathematical Society Colloquium Publications, 55. American Mathematical Society, Providence, RI; Hindustan Book Agency, New Delhi, 2008.

\bibitem{CCentropy}  A.~Connes, C.~Consani,  {\em Characteristic one, entropy and the absolute point}, `` Noncommutative Geometry, Arithmetic, and Related Topics'', the Twenty-First Meeting of the Japan-U.S. Mathematics Institute, Baltimore 2009, JHUP (2012), 75--139.

\bibitem{Cwitt} A.~Connes, {\em The Witt construction in characteristic one  and Quantization}. Noncommutative geometry and global analysis, 83--113, Contemp. Math., 546, Amer. Math. Soc., Providence, RI, 2011.

\bibitem{ccthickening} A.~Connes, C.~Consani, {\em  Universal thickening of the field of real numbers}. Advances in the theory of numbers, 11--74, Fields Inst. Commun., 77, Fields Inst. Res. Math. Sci., Toronto, ON, 2015. 




\bibitem{CC0} A.~Connes, C.~Consani, {\em Schemes over $\F_1$ and zeta functions}, Compositio Mathematica
146 (6), (2010) 1383--1415.

\bibitem{CC0.5} A.~Connes, C.~Consani, {\em From monoids to hyperstructures: in search of an absolute arith-
metic}, in Casimir Force, Casimir Operators and the Riemann Hypothesis, de Gruyter
(2010), 147--198.

\bibitem{CC1} A.~Connes, C.~Consani, {\em The Arithmetic Site}, Comptes Rendus Mathematique Ser. I 352 (2014), 971--975.

\bibitem{CC2} A.~Connes, C.~Consani, {\em Geometry of the Arithmetic Site}. Advances in Mathematics 291 (2016) 274--329.

\bibitem{CC3} A.~Connes, C.~Consani, {\em The Scaling Site}, C.R. Mathematique, Ser. I 354 (2016) 1--6.

\bibitem{CC4} A.~Connes, C.~Consani, {\em Geometry of the Scaling Site}, Selecta Math. New Ser. {\bf 23} no. 3 (2017), 1803-1850.

\bibitem{CC5} A.~Connes, C.~Consani, {Homological algebra in characteristic one},  Preprint (2017), 107 pages.  Available
online at \url{https://arxiv.org/abs/1703.02325}.
%\bibitem{DP} Davey, B. A.; Priestley, H. A. {\em Introduction to Lattices and Order} (second ed.). Cambridge University Press. (2002).

%\bibitem{dejong} DeJong {\em Cohomology of Sheaves}


%\bibitem{dejong1} DeJong {\em Cohomology on sites}
\bibitem{Crh} A.~Connes,  {\em An essay on the Riemann Hypothesis}. In ``Open problems in mathematics", Springer (2016), 
volume edited by Michael Rassias and John Nash.
 
\bibitem{gathmannrr} A.~Gathmann and M.~Kerber. {\em A Riemann-Roch theorem in tropical geometry}. Math. Z.,
259(1):217--230, 2008.

\bibitem{Golan} J.~Golan, {\em Semi-rings and their applications}, Updated and expanded version of The theory of semi-rings, with applications to mathematics and theoretical computer science [Longman Sci. Tech., Harlow, 1992. Kluwer Academic Publishers, Dordrecht, 1999.


\bibitem{gra1} M.~Grandis,  {\em Homological algebra. In strongly non-abelian settings}, World Scientific 2013.

 %\bibitem{Grtohoku} A.~Grothendieck {\em  Sur quelques points d'alg\`ebre homologique}, I. Tohoku Math. J. (2) 9 (1957), no. 2, 119--221.

% \bibitem{MacLane} S.~Mac Lane, {\em Categories for the working mathematician}. Second edition. Graduate Texts in Mathematics, 5. Springer-Verlag, New York, 1998.




%\bibitem{gra} M.~Grandis, {\em A Categorical Approach to Exactness in Algebraic Topology}

 %\bibitem{gra1} M.~Grandis,  {\em Homological algebra}
 
    \bibitem{grmt} A.~Grothendieck {\em Sur une note de Mattuck-Tate} J. reine angew. Math. 200, 208-215 (1958).

\bibitem{Grtohoku} A.~Grothendieck {\em  Sur quelques points d'alg\`ebre homologique}, I. Tohoku Math. J. (2) 9 (1957), no. 2, 119--221

\bibitem{Hart} R.~Hartshorne  {\em Algebraic Geometry}, Graduate Texts in Mathematics 52, Springer-Verlag, New York Heidelberg Berlin 1977.



\bibitem{jessen33}
B.~Jessen, Börge; Über die Nullstellen einer analytischen fastperiodischen Funktion. Eine Verallgemeinerung der Jensenschen Formel. (German) Math. Ann. 108 (1933), no. 1, 485--516.

\bibitem{JT} Jessen, Tornehave  { \em Mean motions and zeros of almost periodic functions}. Acta math. 77, 137--279 (1945).




\bibitem{MacLane} S.~Mac Lane, {\em Categories for the working mathematician}. Second edition. Graduate Texts in Mathematics, 5. Springer-Verlag, New York, 1998.

\bibitem{MM} S.~Mac Lane, I~Moerdijk, {\em Sheaves in geometry and logic. A first introduction to topos theory}. Corrected reprint of the 1992 edition. Universitext. Springer-Verlag, New York, 1994.

\bibitem{Maslov}  V.~Kolokoltsov, V.~P.~Maslov, {\em Idempotent analysis and its applications}.  Mathematics and its Applications, 401. Kluwer Academic Publishers Group, Dordrecht, 1997.

%\bibitem{Li} X.~J.~Li,     {\em The positivity of a sequence of numbers and the Riemann hypothesis},          J. Number Theory 65 (1997), 325--333.

\bibitem{Lit} G.~Litvinov, {\em Tropical Mathematics, Idempotent Analysis, Classical Mechanics and Geometry}. Spectral theory and geometric analysis, 159--186, Contemp. Math., 535, Amer. Math. Soc., Providence, RI, 2011.


\bibitem{Meyer} R.~Meyer, {\em On a representation of the idele class group
related to primes and zeros of $L$-functions}.  Duke Math. J. Vol.127 (2005),  N.3, 519--595.

\bibitem{MZ} G.~Mikhalkin and I.~Zharkov,  {\em Tropical curves, their Jacobians and theta functions}. In Curves and
abelian varieties, volume 465 of Contemp. Math., p 203--230. Amer. Math. Soc., Providence,
RI, 2008.

\bibitem{milne}  J.S.~Milne, {\em Canonical models of Shimura curves},
  manuscript, 2003 (www.jmilne.org).

\bibitem{Neshveyev} M.~Laca, S.~Neshveyev, M.~Trifkovic {\em Bost-Connes systems, Hecke algebras, and induction}.  J. Noncommut. Geom. 7 (2013), no. 2, 525--546.




\bibitem{Robert} A.~Robert,  {\em  A course in p-adic analysis}. Graduate Texts in Mathematics, 198. Springer-Verlag, New York, 2000. 

%\bibitem{Robert} A.~Robert,  {\em  A course in p-adic analysis}. Graduate Texts in Mathematics, 198. Springer-Verlag, New York, 2000. 



%\bibitem{Ronkin}

\bibitem{Rudin} W.~Rudin, {\em Real and complex analysis}. McGraw-Hill, New York, 1966. 



%\bibitem{Tate} J.~Tate  {\em A review of non-archimedean elliptic functions}

%\bibitem{Viro} O.~Viro


\bibitem{vN} J.~von Neumann,  {\em Almost Periodic Functions in a Group I}, Trans. Amer. Math. Soc., 36 no. 3 (1934) pp. 445--492





\bibitem{Weibel} Weibel, Charles A. {\em An introduction to homological algebra}. Cambridge Studies in Advanced Mathematics, 38. Cambridge University Press, Cambridge, 1994. 

\bibitem{weilbasic} A.~Weil, {\em Basic Number Theory}, Reprint of the
second (1973) edition. Classics in Mathematics. Springer-Verlag,
1995.


\bibitem{yoshi} S.~Yoshitomi,  {\em Generators of modules in tropical geometry}. ArXiv Math.AG, 10010448.



\end{thebibliography}
 
 
\end{document}



\begin{thebibliography}{99.}%
% and use \bibitem to create references.
%
% Use the following syntax and markup for your references if 
% the subject of your book is from the field 
% "Mathematics, Physics, Statistics, Computer Science"
%
% Contribution 
\bibitem{science-contrib} Broy, M.: Software engineering --- from auxiliary to key technologies. In: Broy, M., Dener, E. (eds.) Software Pioneers, pp. 10-13. Springer, Heidelberg (2002)
%
% Online Document
\bibitem{science-online} Dod, J.: Effective substances. In: The Dictionary of Substances and Their Effects. Royal Society of Chemistry (1999) Available via DIALOG. \\
\url{http://www.rsc.org/dose/title of subordinate document. Cited 15 Jan 1999}
%
% Monograph
\bibitem{science-mono} Geddes, K.O., Czapor, S.R., Labahn, G.: Algorithms for Computer Algebra. Kluwer, Boston (1992) 
%
% Journal article
\bibitem{science-journal} Hamburger, C.: Quasimonotonicity, regularity and duality for nonlinear systems of partial differential equations. Ann. Mat. Pura. Appl. \textbf{169}, 321--354 (1995)
%
% Journal article by DOI
\bibitem{science-DOI} Slifka, M.K., Whitton, J.L.: Clinical implications of dysregulated cytokine production. J. Mol. Med. (2000) doi: 10.1007/s001090000086 
%
\bigskip

% Use the following (APS) syntax and markup for your references if 
% the subject of your book is from the field 
% "Mathematics, Physics, Statistics, Computer Science"
%
% Online Document
\bibitem{phys-online} J. Dod, in \textit{The Dictionary of Substances and Their Effects}, Royal Society of Chemistry. (Available via DIALOG, 1999), 
\url{http://www.rsc.org/dose/title of subordinate document. Cited 15 Jan 1999}
%
% Monograph
\bibitem{phys-mono} H. Ibach, H. L\"uth, \textit{Solid-State Physics}, 2nd edn. (Springer, New York, 1996), pp. 45-56 
%
% Journal article
\bibitem{phys-journal} S. Preuss, A. Demchuk Jr., M. Stuke, Appl. Phys. A \textbf{61}
%
% Journal article by DOI
\bibitem{phys-DOI} M.K. Slifka, J.L. Whitton, J. Mol. Med., doi: 10.1007/s001090000086
%
% Contribution 
\bibitem{phys-contrib} S.E. Smith, in \textit{Neuromuscular Junction}, ed. by E. Zaimis. Handbook of Experimental Pharmacology, vol 42 (Springer, Heidelberg, 1976), p. 593
%
\bigskip
%
% Use the following syntax and markup for your references if 
% the subject of your book is from the field 
% "Psychology, Social Sciences"
%
%
% Monograph
\bibitem{psysoc-mono} Calfee, R.~C., \& Valencia, R.~R. (1991). \textit{APA guide to preparing manuscripts for journal publication.} Washington, DC: American Psychological Association.
%
% Online Document
\bibitem{psysoc-online} Dod, J. (1999). Effective substances. In: The dictionary of substances and their effects. Royal Society of Chemistry. Available via DIALOG. \\
\url{http://www.rsc.org/dose/Effective substances.} Cited 15 Jan 1999.
%
% Journal article
\bibitem{psysoc-journal} Harris, M., Karper, E., Stacks, G., Hoffman, D., DeNiro, R., Cruz, P., et al. (2001). Writing labs and the Hollywood connection. \textit{J Film} Writing, 44(3), 213--245.
%
% Contribution 
\bibitem{psysoc-contrib} O'Neil, J.~M., \& Egan, J. (1992). Men's and women's gender role journeys: Metaphor for healing, transition, and transformation. In B.~R. Wainrig (Ed.), \textit{Gender issues across the life cycle} (pp. 107--123). New York: Springer.
%
% Journal article by DOI
\bibitem{psysoc-DOI}Kreger, M., Brindis, C.D., Manuel, D.M., Sassoubre, L. (2007). Lessons learned in systems change initiatives: benchmarks and indicators. \textit{American Journal of Community Psychology}, doi: 10.1007/s10464-007-9108-14.
%
%
% Use the following syntax and markup for your references if 
% the subject of your book is from the field 
% "Humanities, Linguistics, Philosophy"
%
\bigskip
%
% Journal article
\bibitem{humlinphil-journal} Alber John, Daniel C. O'Connell, and Sabine Kowal. 2002. Personal perspective in TV interviews. \textit{Pragmatics} 12:257--271
%
% Contribution 
\bibitem{humlinphil-contrib} Cameron, Deborah. 1997. Theoretical debates in feminist linguistics: Questions of sex and gender. In \textit{Gender and discourse}, ed. Ruth Wodak, 99--119. London: Sage Publications.
%
% Monograph
\bibitem{humlinphil-mono} Cameron, Deborah. 1985. \textit{Feminism and linguistic theory.} New York: St. Martin's Press.
%
% Online Document
\bibitem{humlinphil-online} Dod, Jake. 1999. Effective substances. In: The dictionary of substances and their effects. Royal Society of Chemistry. Available via DIALOG. \\
http://www.rsc.org/dose/title of subordinate document. Cited 15 Jan 1999
%
% Journal article by DOI
\bibitem{humlinphil-DOI} Suleiman, Camelia, Daniel C. O�Connell, and Sabine Kowal. 2002. `If you and I, if we, in this later day, lose that sacred fire...�': Perspective in political interviews. \textit{Journal of Psycholinguistic Research}. doi: 10.1023/A:1015592129296.
%
%
%
\bigskip
%
%
% Use the following syntax and markup for your references if 
% the subject of your book is from the field 
% "Computer Science, Economics, Engineering, Geosciences, Life Sciences"
%
%
% Contribution 
\bibitem{basic-contrib} Brown B, Aaron M (2001) The politics of nature. In: Smith J (ed) The rise of modern genomics, 3rd edn. Wiley, New York 
%
% Online Document
\bibitem{basic-online} Dod J (1999) Effective Substances. In: The dictionary of substances and their effects. Royal Society of Chemistry. Available via DIALOG. \\
\url{http://www.rsc.org/dose/title of subordinate document. Cited 15 Jan 1999}
%
% Journal article by DOI
\bibitem{basic-DOI} Slifka MK, Whitton JL (2000) Clinical implications of dysregulated cytokine production. J Mol Med, doi: 10.1007/s001090000086
%
% Journal article
\bibitem{basic-journal} Smith J, Jones M Jr, Houghton L et al (1999) Future of health insurance. N Engl J Med 965:325--329
%
% Monograph
\bibitem{basic-mono} South J, Blass B (2001) The future of modern genomics. Blackwell, London 
%
\end{thebibliography}
